% Theory of valuations: Plan B project for the master's degree at Central Michigan University.
% The website page notes/projects/valuations.html and the PDF are both made from this file.
% The two Python programs are in valuations-code/; the web page runs them in the browser (notes/projects/padic-tools.js).
\documentclass[11pt]{article}
\usepackage{pdenotes}
\usepackage{listings}
\lstset{basicstyle=\ttfamily\footnotesize,breaklines=true,columns=fixed,keepspaces=true,basewidth=0.5em,frame=single,framerule=0.3pt,
  rulecolor=\color{black!25},backgroundcolor=\color{boxbg},showstringspaces=false}
\hypersetup{pdftitle={Theory of Valuations},pdfauthor={Amer Rawashdeh}}
\pagestyle{plain}
\title{Theory of Valuations}

\begin{document}
\begin{pdfonly}
\begin{center}
{\LARGE\bfseries Theory of Valuations\par}\medskip
{\large Amer Rawashdeh\par}
\end{center}
\end{pdfonly}

\begin{byline}
Plan B paper for the master's degree in mathematics, Central Michigan University.
\end{byline}

\begin{webonly}
The two Python programs of this project run on this page: see \href{\#programs}{the last section}.
Also as a \href{../../latex/projects/valuations.pdf}{printable PDF}.
\end{webonly}

\paragraph{Introduction.}
The valuation of a field/ring is a way to study abstract structures using analysis methods. For example, if we have a valued field/ring we can talk about completeness or study sequences. Also, in this paper I will talk about the $p$-adic numbers for $p$ a prime, which we can think about as a base-$p$ number system (binary when $p=2$). Also, we can ask interesting questions in the field of $p$-adic numbers like: is $\sqrt{2}$ a $p$-adic number? i.e.\ for example, does the equation $x^2-2=0$ have a solution in the $7$-adic or the $5$-adic numbers?

\paragraph{Abstract.}
In this paper I mostly used the book Algebraic Number Theory by Neukirch [1]. I talked about the $p$-adic numbers, gave some examples and wrote two codes: one code to convert a rational number into its $p$-adic representation, and the second code to check if a number has a square root in the $p$-adic numbers for a certain $p$. In the next sections I talked about the theory of valuations, the $p$-adic valuation, the field of rational functions and its completion with respect to a nonarchimedean valuation (formal Laurent series), and Puiseux series.

\section{$p$-adic numbers}

Let $p$ be a prime number. The \index{p-adic!Integers}\textit{$p$-adic integers} $\mathbb{Z}_p$ is the set with the elements in the form
\begin{equation}\label{eq:val-zp}
\mathbb{Z}_p = \left\{ \sum_{i=0}^{\infty} a_i p^i : 0 \leq a_i < p \right\}.
\end{equation}
$\mathbb{Z}_p$ with two operations, addition and multiplication (with carrying, as for numbers written in base $p$), forms a ring. And its field of fractions is the set of \index{p-adic!Numbers}\textit{$p$-adic numbers} $\mathbb{Q}_p$:
\begin{equation}\label{eq:val-qp}
\mathbb{Q}_p = \left\{ \sum_{i=m}^\infty a_i p^i : m \in \mathbb{Z},\ 0 \leq a_i < p \right\}.
\end{equation}

We can write the elements in $\mathbb{Q}_p$ by the coefficients of the $p$ powers as $a_m a_{m+1} a_{m+2} \cdots$, and we use a comma after the $a_0$ coefficient (the coefficient with no $p$ power). There is a 1-1 map (embedding) from the set of rational numbers $\mathbb{Q}$ into $\mathbb{Q}_p$: a rational number corresponds to the element of $\mathbb{Q}_p$ whose sum in \eqref{eq:val-qp} has that value, when we compute the sum without concern about the convergence (as in the next example).

\begin{exbox}{Example 1.1}
Find the rational number which corresponds to the $7$-adic number $3,545454\cdots$
\[
\begin{aligned}
3,54545454\cdots &= 3 + 5\cdot 7 + 4\cdot 7^2 + 5\cdot 7^3 + 4\cdot 7^4 + \cdots \\
&= 3 + (5\cdot 7 + 4\cdot 7^2)(1 + 7^2 + 7^4 + \cdots) \\
&= 3 + (5\cdot 7 + 4\cdot 7^2)\left(\frac{1}{1-7^2}\right) \\
&= 3 + \frac{35+196}{-48} = \frac{-144+35+196}{-48} = \frac{-87}{48} = \frac{-29}{16}.
\end{aligned}
\]
\end{exbox}

Also, if we want to convert from the rationals to the $p$-adics, we do that coefficient by coefficient, in the following way.

Let $c \in \mathbb{Q}$. Then we factor the $p$ powers from $c$, so we have $c = p^m b_m$ where $p \nmid b_m$ (that is, $p$ divides neither the numerator nor the denominator of $b_m$). Now for $b_m$ we find an integer $0 \leq a_m < p$ such that $p \mid b_m - a_m$, and we write $b_m = a_m + p\cdot b_{m+1}$. This $a_m$ is the coefficient of $p^m$. Then we do the same for $b_{m+1}$ to find $a_{m+1}$, the coefficient of $p^{m+1}$, and we keep going.

\begin{exbox}{Example 1.2}
What does $\frac{110}{7}$ correspond to in the $11$-adics? First we factor the $11$ powers:
\[
\frac{110}{7} = 11 \cdot \frac{10}{7}.
\]
Here we will begin with finding $a_1$, because the fraction has $p^1$ as a factor. So we want to find $0 \leq a_1 < 11$ such that $\frac{10}{7} - a_1$ has $11$ as a factor. The solution is $3$; to see this,
\[
\frac{10}{7} - 3 = \frac{10-21}{7} = 11\cdot \frac{-1}{7}.
\]
So $a_1 = 3$. Now we want to find $a_2$ such that $\frac{-1}{7} - a_2$ has $11$ as a factor; again it will be $3$:
\[
\frac{-1}{7} - 3 = \frac{-1-21}{7} = 11 \cdot \frac{-2}{7}.
\]
So $a_2 = 3$. Now we want to find $a_3$ such that $\frac{-2}{7} - a_3$ has $11$ as a factor. The answer is $6$:
\[
\frac{-2}{7} - 6 = \frac{-44}{7} = 11 \cdot \frac{-4}{7}.
\]
So $a_3 = 6$. And then $a_4 = 1$:
\[
\frac{-4}{7} - 1 = \frac{-4-7}{7} = 11 \cdot \frac{-1}{7}.
\]
Then it will be periodic, because we get $\frac{-1}{7}$, which we already had. Hence
\[
\frac{110}{7} = 0,3\; 3\; 6\; 1\; 3\; 6\; 1 \cdots
\]
We can check that it is correct in two ways. First, we use the geometric sum:
\[
0,3361361\cdots = 3\cdot 11 + (3\cdot 11^2 + 6\cdot 11^3 + 11^4)(1 + 11^3 + \cdots) = 33 - \frac{3\cdot 11^2 + 6\cdot 11^3 + 11^4}{11^3 - 1} = \frac{110}{7}.
\]
Second, if we multiply $\frac{110}{7}$ by $7$ we get $110$, which as an $11$-adic is $0,10$ (here $10$ is one coefficient, since $110 = 10\cdot 11$):
\[
7\,(0,3361361\cdots) = 0,10\; 0\; 0\; 0 \cdots
\]
Multiplying $7$ with $3$ we get $21 = 10 + 1\cdot 11$, so the carried number is $1$; then $7\cdot 3 + 1 = 22 = 0 + 2\cdot 11$, so the carried number is $2$; then $7\cdot 6 + 2 = 44 = 0 + 4\cdot 11$, so the carried number is $4$; then $7\cdot 1 + 4 = 11 = 0 + 1\cdot 11$, and so on.
\end{exbox}

There is an interesting exercise in [1, Exercise 1, p.~106]:

\begin{thmbox}{Proposition 1.3}
A $p$-adic number $a = \sum_{i=m}^\infty a_i p^i$ is a rational number if and only if the sequence of coefficients $a_m, a_{m+1}, \cdots$ is periodic after a finite number of terms.
\end{thmbox}

\begin{proof}
$(\Rightarrow)$ Suppose that $a \in \mathbb{Q}_p$ is a rational number, say $\frac{s}{t}$. We want to show the $p$-adic representation will be periodic after a finite number of terms. First, we factor the $p$ powers so that
\[
\frac{s}{t} = p^m \cdot \frac{s'}{t'}, \quad \text{with } p \nmid s't' \text{ and } t' > 0 .
\]
We use Euclid division $s' = q\cdot t' + r$ with $q \in \mathbb{Z}$ and $-t' \leq r \leq -1$ (we choose $q$ so that the remainder is negative), so we have $\frac{s'}{t'} = q + \frac{r}{t'}$. Since $t'$ is relatively prime with $p$, we can use Euler's theorem (the generalization of Fermat's little theorem) to get
\[
p^{\varphi(t')} \equiv 1 \pmod{t'},
\]
where $\varphi(t')$ is the cardinality of the set $A = \{ x : 1 \leq x \leq t' \ \&\ \gcd(x,t') = 1 \}$. Thus $p^{\varphi(t')} - 1 = kt'$ with $k \geq 1$. Let $\varphi(t') = N$ and let $x = -kr$. Then
\[
\frac{r}{t'} = \frac{kr}{kt'} = \frac{-x}{p^N - 1} = \frac{x}{1 - p^N} = x \underbrace{(1 + p^N + p^{2N} + \cdots)}_{1/(1-p^N)} .
\]
Since $x \in \mathbb{N}$, we write its $p$-adic expansion, and $a_N = a_{N+1} = \cdots = 0$ in its expansion because $-t' \leq r \leq -1 \Rightarrow 0 < x = -kr \leq kt' = p^N - 1 < p^N$. Suppose $x = a_0 + a_1 p + \cdots + a_{N-1} p^{N-1}$. Then we get that
\[
\frac{r}{t'} = a_0 + a_1 p + \cdots + a_{N-1} p^{N-1} + a_0 p^N + a_1 p^{1+N} + \cdots + a_{N-1} p^{2N-1} + \cdots
\]
which is periodic. Thus
\[
\frac{s}{t} = p^m \cdot \left( q + a_0 + a_1 p + \cdots + a_{N-1} p^{N-1} + a_0 p^N + a_1 p^{1+N} + \cdots \right).
\]
Multiplying by $p^m$ only shifts the coefficients. I will show that adding an integer does not affect the periodicity of a $p$-adic number. Since any integer $q$ is a sum of $1$'s or of $-1$'s, if adding $1$ and adding $-1$ do not affect the periodicity of a $p$-adic number, then adding an integer will not. Let $a_0 + a_1 p + \cdots$ be the $p$-adic integer. If $a_0 < p-1$, then adding $1$ will just affect the first coefficient. If $a_0 = p-1$, then $a_0$ becomes $0$ and a $1$ will carry on to the next coefficient, and so on. So adding $1$ changes only finitely many coefficients, unless all coefficients from some point on are $p-1$; then those all become $0$. In both cases the result is periodic after a finite number of terms. Adding $-1$ is the same with borrowing instead of carrying: it changes only finitely many coefficients, unless all coefficients from some point on are $0$, and then those all become $p-1$.

$(\Leftarrow)$ Let $a = \sum_{i=m}^\infty a_i p^i$; multiplying by $p^{-m}$ we can assume $m = 0$. Suppose it is periodic after $N$ terms, and suppose the length of the string which repeats is $k$. Then
\[
\begin{aligned}
a &= a_0 + \cdots + a_{N-1}p^{N-1} + (a_N p^N + \cdots + a_{N+k-1}p^{N+k-1})(1 + p^k + p^{2k} + \cdots) \\
&= a_0 + \cdots + a_{N-1}p^{N-1} + \frac{a_N p^N + \cdots + a_{N+k-1}p^{N+k-1}}{1 - p^k},
\end{aligned}
\]
which is a rational number.
\end{proof}

\paragraph{Another presentation of the $p$-adics.} There is another way to write the $p$-adic integers: instead of writing the coefficients $a_0, a_1 a_2 \cdots$ we can write the partial sums of $\sum_{i=0}^\infty a_i p^i$. We define
\[
\begin{aligned}
s_0 &= a_0, \\
s_1 &= a_0 + a_1\cdot p, \\
s_2 &= a_0 + a_1\cdot p + a_2\cdot p^2 .
\end{aligned}
\]
In general it will be
\[
s_n = \sum_{i=0}^{n} a_i p^i .
\]
Define the \index{Inverse limit}\textbf{inverse (projective) limit} of the rings $\mathbb{Z}/p^n\mathbb{Z}$ to be the set
\[
\varprojlim \mathbb{Z}/p^n\mathbb{Z} = \left\{ \prod_{i=0}^\infty (s_i) : s_i \in \mathbb{Z}/p^{i+1}\mathbb{Z} \ \&\ s_{i+1} \equiv s_i \pmod{p^{i+1}} \right\}.
\]
Now we need to show that the two representations of the $p$-adics are the same. We define the map
\[
\phi : \mathbb{Z}_p \rightarrow \varprojlim \mathbb{Z}/p^n\mathbb{Z} \quad\text{by}\quad \phi(a_0, a_1 a_2 \cdots) = \prod_{j=0}^\infty \left( \sum_{i=0}^j a_i p^i \right).
\]
It is a bijection, with the inverse
\[
\phi^{-1}\left( \prod_{i=0}^\infty (s_i) \right) = \left( s_0, \frac{s_1 - s_0}{p}, \frac{s_2 - s_1}{p^2}, \frac{s_3 - s_2}{p^3}, \cdots \right),
\]
where each $s_i$ is taken as an integer with $0 \leq s_i < p^{i+1}$. And this representation is useful in proving the following proposition.

\begin{thmbox}{Proposition 1.4}
Let $Q(x)$ be a polynomial with integer coefficients. Then the equation $Q(x) = 0$ has a solution in $\mathbb{Z}_p$ iff $Q(x) \equiv 0 \pmod{p^n}$ has a solution for every $n$.
\end{thmbox}

\begin{proof}
$(\Rightarrow)$ Suppose that the equation $Q(x) = 0$ has a solution in the $p$-adics, i.e.\ there is $x_p \in \mathbb{Z}_p$ such that $Q(x_p) = 0$. Then, writing $s_{n-1}$ for the partial sum of $x_p$ up to $p^{n-1}$, we have $x_p \equiv s_{n-1} \pmod{p^n}$, so clearly from the mod properties $Q(s_{n-1}) \equiv Q(x_p) = 0 \pmod{p^n}$ for any $n \in \mathbb{N}$.

$(\Leftarrow)$ Now suppose that $Q(x) \equiv 0 \pmod{p^n}$ has a solution for any $n \in \mathbb{N}$. Note that a solution mod $p^{N}$ is also a solution mod $p^n$ for every $n \leq N$. Call a solution $t$ mod $p^n$ \emph{good} if for every $N \geq n$ there is a solution mod $p^N$ which is $\equiv t \pmod{p^n}$.

First, there is a good solution mod $p$. Indeed, every solution mod $p^N$ reduces mod $p$ to one of the finitely many elements of $\mathbb{Z}/p\mathbb{Z}$, so one element $s_0$ is the reduction of solutions mod $p^N$ for infinitely many $N$, and hence for all $N$ (a solution mod a higher power reduces to a solution mod a lower power).

Next, if $s_n$ is a good solution mod $p^{n+1}$, then one of its $p$ lifts $s_n + c\,p^{n+1}$, $0 \leq c < p$, is a good solution mod $p^{n+2}$, by the same argument: the solutions mod $p^N$ ($N \geq n+2$) which are $\equiv s_n \pmod{p^{n+1}}$ reduce mod $p^{n+2}$ to one of these $p$ lifts, so one of them is used for infinitely many $N$, and hence for all $N$. Call it $s_{n+1}$.

So we get $s_0, s_1, s_2, \cdots$ with $s_i \in \mathbb{Z}/p^{i+1}\mathbb{Z}$, $s_{i+1} \equiv s_i \pmod{p^{i+1}}$ and $Q(s_i) \equiv 0 \pmod{p^{i+1}}$. This is an element of $\varprojlim \mathbb{Z}/p^n\mathbb{Z}$, so by the bijection $\phi$ it is a $p$-adic integer $x_p = \phi^{-1}\left(\prod (s_i)\right)$. Then $Q(x_p) \equiv Q(s_i) \equiv 0 \pmod{p^{i+1}}$ for every $i$, i.e.\ all the coefficients of $Q(x_p)$ are $0$, so $Q(x_p) = 0$.
\end{proof}

\begin{exbox}{Example 1.5}
We will show that $\sqrt{2}$ is in the $7$-adics but not in the $11$-adics. To prove that $\sqrt{2}$ is in the $7$-adics, it is enough, by the previous proposition, to show there is a solution to the following congruence:
\[
x^2 \equiv 2 \pmod{7^n}, \quad \forall n \in \mathbb{N}.
\]
I will use induction. For $n = 1$ it is easy to see that $3^2 \equiv 2 \pmod 7$. Now suppose that
\[
x^2 \equiv 2 \pmod{7^n}
\]
is solvable, and let one of the solutions be
\[
x_n = \sum_{i=0}^{n-1} a_i \cdot 7^{i}, \qquad a_0 = 3 .
\]
Now I will show there is $a_n \in \{0, 1, 2, \cdots, 6\}$ such that
\[
\underbrace{(x_n + a_n \cdot 7^{n})^2}_{x_{n+1}^2} \equiv 2 \pmod{7^{n+1}} .
\]
We know that $x_n^2 \equiv 2 \pmod{7^n}$, which means that $7^n \mid x_n^2 - 2$. I will use that to solve the congruence:
\[
\begin{aligned}
(x_n + a_n \cdot 7^{n})^2 &\equiv 2 \pmod{7^{n+1}} \\
x_n^2 - 2 + 2x_n a_n \cdot 7^{n} + {\color{red} a_n^2\, 7^{2n}} &\equiv 0 \pmod{7^{n+1}} \\
7^n \left( \frac{x_n^2 - 2}{7^n} + 2x_n a_n \right) &\equiv 0 \pmod{7^{n+1}}
\end{aligned}
\]
(the red term is divisible by $7^{n+1}$). Thus
\[
\frac{x_n^2 - 2}{7^n} + 2x_n a_n \equiv 0 \pmod 7 .
\]
We will simplify it (all congruences below are mod $7$):
\[
\begin{aligned}
\frac{x_n^2 - 2}{7^n} + 2x_n a_n &\equiv 0 && (1) \\
2x_n a_n &\equiv -\frac{x_n^2 - 2}{7^n} && (2) \\
6\cdot a_n &\equiv 6\cdot \frac{x_n^2 - 2}{7^n} && (3) \\
a_n &\equiv \frac{x_n^2 - 2}{7^n} && (4) \\
a_n &\equiv \frac{(x_{n-1} + a_{n-1}\cdot 7^{n-1})^2 - 2}{7^n} && (5) \\
a_n &\equiv \frac{x_{n-1}^2 + 2a_{n-1}x_{n-1}\cdot 7^{n-1} + a_{n-1}^2\cdot 7^{2n-2} - 2}{7^n} && (6) \\
a_n &\equiv \frac{x_{n-1}^2 - 2 + 2a_{n-1}x_{n-1}\cdot 7^{n-1}}{7^n} + a_{n-1}^2\cdot 7^{n-2} && (7) \\
a_n &\equiv \frac{x_{n-1}^2 - 2 + 2a_{n-1}x_{n-1}\cdot 7^{n-1}}{7^n} && (8)
\end{aligned}
\]
Here are the details of what I did:
\begin{itemize}
\item from (2) to (3) I used the fact that $x_n \equiv 3 \pmod 7$, so $2x_n \equiv 6 \pmod 7$, and $-1 \equiv 6 \pmod 7$;
\item from (3) to (4) I multiplied by $6$, because it is its own inverse: $6\cdot 6 = 36 \equiv 1 \pmod 7$;
\item from (4) to (5) using that $x_n = x_{n-1} + a_{n-1}\cdot 7^{n-1}$;
\item from (7) to (8): since we are taking mod $7$, when $n \geq 3$ we get $a_{n-1}^2 \cdot 7^{n-2} \equiv 0$.
\end{itemize}
Let's compute the first four terms. We know that $a_0 = 3$, and to find $a_1$ I will use (4):
\[
a_1 \equiv \frac{3^2 - 2}{7} \equiv 1 .
\]
Now to find $a_2$ I will use (4) again, with $x_2 = 3 + 7 = 10$:
\[
a_2 \equiv \frac{10^2 - 2}{7^2} \equiv 2 .
\]
We have $x_3 = 3 + 7 + 2\cdot 7^2$; to find $a_3$ I will use (8):
\[
a_3 \equiv \frac{x_2^2 - 2 + 2a_2 x_2\cdot 7^2}{7^3} \equiv \frac{10^2 - 2 + 2(2)(10)\cdot 7^2}{7^3} = \frac{2058}{343} = 6 .
\]
So $x_4 = 3 + 7 + 2\cdot 7^2 + 6\cdot 7^3$.

To show that $\sqrt{2}$ is not in the $11$-adics, we look at the elements of $\mathbb{Z}/11\mathbb{Z}$, which are $\{0, 1, 2, \cdots, 10\}$; their squares are $\{0, 1, 4, 9, 5, 3\}$. Hence
\[
x^2 \equiv 2 \pmod{11}
\]
has no solution, so $\sqrt{2}$ is not an element of the $11$-adic integers. (It is not in $\mathbb{Q}_{11}$ either: if $x = 11^m u$ with $u \in \mathbb{Z}_{11}$, $11 \nmid u$, and $x^2 = 2$, then $11^{2m}u^2 = 2$ forces $m = 0$.)
\end{exbox}

What I did to show that $\sqrt{2}$ is in the $7$-adics and to evaluate the first $4$ terms has a generalization in the following lemma.

\begin{thmbox}{Hensel's Lemma (first version)}
\index{Hensel's Lemma}If $Q(x)$ is a polynomial with integer coefficients and $Q(x) \equiv 0 \pmod{p^n}$ has a solution, say $t \in \mathbb{Z}/p^n\mathbb{Z}$, then this solution can be lifted to a solution $\pmod{p^{n+1}}$ if $Q'(t) \not\equiv 0 \pmod p$.
\end{thmbox}

\begin{proof}
Suppose $Q(x) = \sum_{i=0}^m a_i x^i$ and suppose
\[
Q(t) = \sum_{i=0}^m a_i t^i \equiv 0 \pmod{p^n} .
\]
We need to look for $0 \leq c < p$ such that $t + cp^n$ satisfies $Q(t + cp^n) \equiv 0 \pmod{p^{n+1}}$. Then
\[
\begin{aligned}
\sum_{i=0}^m a_i (t + cp^n)^i &\equiv 0 \pmod{p^{n+1}} \\
a_0 + \sum_{i=1}^m a_i \Big( t^i + i t^{i-1}(cp^n) + \underbrace{\tfrac{i(i-1)}{2} t^{i-2}(cp^n)^2 + \cdots + (cp^n)^i}_{\equiv 0 \pmod{p^{n+1}}} \Big) &\equiv 0 \pmod{p^{n+1}} \\
a_0 + \sum_{i=1}^m a_i \big( t^i + i t^{i-1}(cp^n) \big) &\equiv 0 \pmod{p^{n+1}} \\
Q(t) + cp^n Q'(t) &\equiv 0 \pmod{p^{n+1}} .
\end{aligned}
\]
Dividing by $p^n$, this is $\frac{Q(t)}{p^n} + c\,Q'(t) \equiv 0 \pmod p$, and this last congruence has a solution $c$ if $Q'(t) \not\equiv 0 \pmod p$, namely $c \equiv -\frac{Q(t)}{p^n}\, Q'(t)^{-1} \pmod p$.
\end{proof}

\section{Valuations}

A \index{Valuation}\textit{valuation} on a field $\mathbb{F}$ is a function $|\cdot| : \mathbb{F} \rightarrow \mathbb{R}$ satisfying
\begin{itemize}
\item $|x| \geq 0$ for any $x \in \mathbb{F}$,
\item $|x| = 0$ iff $x = 0$,
\item $|xy| = |x||y|$,
\item $|x + y| \leq |x| + |y|$.
\end{itemize}
From the definition of the valuation, if $1$ is the multiplicative identity of $\mathbb{F}$, then $|1|\cdot|1| = |1\cdot 1| = |1|$, so $|1| = 1$.

A valuation is called \textit{nonarchimedean}\index{Valuation!Nonarchimedean} if there is $M \in \mathbb{R}$ such that $|n| \leq M$ for all $n \in \mathbb{N}$; otherwise it is called \textit{archimedean}\index{Valuation!Archimedean}.

\begin{thmbox}{Proposition 2.1}
A valuation is nonarchimedean iff it satisfies the \index{Strong triangle inequality}\textit{strong triangle inequality}, i.e.\ $|x + y| \leq \max\{|x|, |y|\}$.
\end{thmbox}

\begin{proof}
$(\Rightarrow)$ Suppose that a valuation $|\cdot|$ on the field $\mathbb{F}$ is nonarchimedean, i.e.\ there is $M \in \mathbb{R}$ such that $|n| \leq M$ for all $n \in \mathbb{N}$. Let $x, y \in \mathbb{F}$ be nonzero elements and suppose that $\max\{|x|, |y|\} = |y|$. Then
\[
\begin{aligned}
|x + y|^n &= |(x + y)^n| \\
&= \left| \sum_{i=0}^n \binom{n}{i} x^i y^{n-i} \right| \\
&\leq \sum_{i=0}^n \left| \binom{n}{i} \right| |y|^n \\
&\leq M\, |y|^n (n + 1) .
\end{aligned}
\]
Hence $|x + y| \leq M^{\frac1n} |y| (n + 1)^{\frac1n}$; if $n \rightarrow \infty$ then we get $|x + y| \leq |y| = \max\{|x|, |y|\}$.

$(\Leftarrow)$ Suppose that a valuation on a field $\mathbb{F}$ satisfies $|x + y| \leq \max\{|x|, |y|\}$ for any $x, y \in \mathbb{F}$. Let $n \in \mathbb{N}$; then
\[
|n| = |(n - 1) + 1| \leq \max\{|n - 1|, |1|\} .
\]
Thus, by induction, $|n| \leq |1|$ for any $n \in \mathbb{N}$, so we set $M = |1| = 1$.
\end{proof}

\begin{exbox}{Example 2.2}
We can define a nonarchimedean valuation on $\mathbb{Q}$. Let $p$ be a prime. If $x \in \mathbb{Q}$, $x \neq 0$, is such that $x = p^m \cdot \frac{a}{b}$ with $p \nmid ab$, then we define the map $v_p : \mathbb{Q} \rightarrow \mathbb{Z} \cup \{\infty\}$ as $v_p(x) = m$. Also, we let $v_p(0) = \infty$.

This $v_p$ is called the \index{Valuation!Exponential}\textit{exponential valuation}, and it satisfies the following properties:
\begin{itemize}
\item $v_p(x) = \infty$ iff $x = 0$,
\item $v_p(xy) = v_p(x) + v_p(y)$,
\item $v_p(x + y) \geq \min\{v_p(x), v_p(y)\}$.
\end{itemize}
To prove the second, suppose $x = p^m u$, $y = p^n w$ where $p$ divides no numerator or denominator of $u, w$; then $xy = p^{m+n}uw$, so clearly $v_p(xy) = m + n = v_p(x) + v_p(y)$. To prove the third property, let again $x = p^m u$, $y = p^n w$, and suppose that $m < n$, i.e.\ $v_p(x) < v_p(y)$. Then
\[
v_p(x + y) = v_p\big(p^m(u + p^{n-m}w)\big) = m \geq \min\{v_p(x), v_p(y)\}.
\]
If $m = n$, then $x + y = p^m(u + w)$, and $u + w$ has no $p$ in its denominator, so $v_p(x + y) \geq m = \min\{v_p(x), v_p(y)\}$. (If $x$ or $y$ is $0$ the property is clear.)

Then we define the nonarchimedean valuation $|\cdot|_p : \mathbb{Q} \rightarrow \mathbb{R}$ by
\[
|x|_p = p^{-v_p(x)} \qquad (\text{with } |0|_p = p^{-\infty} = 0).
\]
From this definition we can observe the following properties:
\begin{itemize}
\item $|x|_p \geq 0$ for all $x \in \mathbb{Q}$,
\item $|x|_p = 0$ iff $x = 0$,
\item $|xy|_p = |x|_p |y|_p$,
\item $|x + y|_p \leq \max\{|x|_p, |y|_p\}$.
\end{itemize}
The first is clear, since $|x|_p = p^{-v_p(x)}$, which is always positive or $0$. The second property is true from the fact $\lim_{t \rightarrow \infty} p^{-t} = 0$, and $v_p(x) = \infty$ only for $x = 0$. The third: $|xy|_p = p^{-v_p(xy)} = p^{-v_p(x) - v_p(y)} = p^{-v_p(x)} p^{-v_p(y)} = |x|_p |y|_p$. The fourth one: $|x + y|_p = p^{-v_p(x + y)} \leq p^{-\min\{v_p(x), v_p(y)\}} = \max\{|x|_p, |y|_p\}$. So $|\cdot|_p$ defines a valuation on $\mathbb{Q}$.
\end{exbox}

\begin{exbox}{Example 2.3}
Consider the $3$-adic valuation on $\mathbb{Q}$. For any integer $n \in \mathbb{Z}$ we get that $|n|_3 \leq 1$. Let $A = \{\frac{a}{b} : 1 < |\frac{a}{b}|_3 < 4\}$; this set contains all rational numbers $\frac{a}{b}$ (in lowest terms) such that $3 \mid b$ but $3^2 \nmid b$, since these are exactly the ones with $|\frac{a}{b}|_3 = 3$. The $3$-adic valuation defines a metric on the rational numbers; for instance $d(\frac13, 10) = |\frac13 - 10|_3 = |\frac{-29}{3}|_3 = 3$. Also, for any $a, b \in \mathbb{Z}$ with $3 \nmid b$ we have $|\frac13 - \frac{a}{b}|_3 = 3$. Suppose $b = 3^r c$ where $r \geq 2$, $3 \nmid c$ and $3 \nmid a$. Then
\[
d\left(\frac13, \frac{a}{b}\right) = \left| \frac13 - \frac{a}{b} \right|_3 = \left| \frac{3^{r-1}c - a}{3^r c} \right|_3 = 3^r ,
\]
since $3 \nmid 3^{r-1}c - a$. In general, suppose $r_1, r_2 \in \mathbb{Q}$ and $r_1 = p^{e_1}q_1$, $r_2 = p^{e_2}q_2$ with $v_p(q_1) = v_p(q_2) = 0$ and $e_1 < e_2$. Then
\[
d_p(r_1, r_2) = |p^{e_1}q_1 - p^{e_2}q_2|_p = |p^{e_1}(q_1 - p^{e_2 - e_1}q_2)|_p = p^{-e_1} = |r_1|_p .
\]
\end{exbox}

\begin{thmbox}{Proposition 2.4}
For any nonzero rational number $a \in \mathbb{Q}$ we have
\[
\prod_{p \text{ prime}} |a|_p = \frac{1}{\|a\|},
\]
where $\|a\|$ is the absolute value of $a$.
\end{thmbox}

\begin{proof}
We write the factorization of $a$ as $a = \operatorname{sign}(a)\, p_1^{e_1} p_2^{e_2} \cdots p_n^{e_n}$ with $e_i \in \mathbb{Z}$; then $|a|_{p_i} = \frac{1}{p_i^{e_i}}$, and $|a|_p = 1$ for every other prime $p$. Thus
\[
\prod_{p \text{ prime}} |a|_p = \frac{1}{p_1^{e_1} p_2^{e_2} \cdots p_n^{e_n}} = \frac{1}{\|a\|} .
\]
\end{proof}

We say two valuations $|\cdot|_a$, $|\cdot|_b$ on $\mathbb{F}$ are \textit{equivalent}\index{Valuation!Equivalent} if there exists $s \in \mathbb{R}^+$ such that $|x|_a = |x|_b^s$ for all $x \in \mathbb{F}$.

\begin{thmbox}{Proposition 2.5}
Two valuations are equivalent if and only if they induce the same topology. (The following is a recap of what is in [2].)
\end{thmbox}

\begin{proof}
Let $|\cdot|_1, |\cdot|_2$ be two valuations on the same field $\mathbb{K}$. Suppose the two valuations are equivalent, i.e.\ there is $s \in \mathbb{R}^+$ such that $|x|_1 = |x|_2^s$ for all $x \in \mathbb{K}$. Then the balls
\[
B_1(a, r^s) = \{ x \in \mathbb{K} : |x - a|_1 < r^s \} = B_2(a, r) = \{ x \in \mathbb{K} : |x - a|_2 < r \}
\]
are the same. So the two topologies induced by $|\cdot|_1$ and $|\cdot|_2$ are the same.

Conversely, suppose the two topologies $\tau_1, \tau_2$ induced respectively by $|\cdot|_1$ and $|\cdot|_2$ are the same. First, we will show that $|x_n|_1 \rightarrow 0 \Longleftrightarrow |x_n|_2 \rightarrow 0$ for a sequence $x_n \in \mathbb{K}$. By symmetry it is enough to show that $|x_n|_1 \rightarrow 0 \Rightarrow |x_n|_2 \rightarrow 0$. Let $\epsilon > 0$. Then $B_2(0, \epsilon)$ is open in $\tau_2 = \tau_1$, and $0 \in B_2(0, \epsilon)$, so there is $\epsilon' > 0$ such that
\[
B_1(0, \epsilon') \subset B_2(0, \epsilon) .
\]
But $|x_n|_1 \rightarrow 0$, so there is $N$ such that $x_n \in B_1(0, \epsilon') \subset B_2(0, \epsilon)$ for all $n > N$. So $|x_n|_2 \rightarrow 0$.

Second, $|x|_1 < 1 \Longleftrightarrow |x|_2 < 1$. Indeed, $|x|_1 < 1$ iff $|x^n|_1 = |x|_1^n \rightarrow 0$ (if $|x|_1 \geq 1$ then $|x^n|_1 \geq 1$ for all $n$), iff $|x^n|_2 \rightarrow 0$ by the first step, iff $|x|_2 < 1$. Applying this to $x^{-1}$, we also get $|x|_1 > 1 \Longleftrightarrow |x|_2 > 1$.

If $|\cdot|_1$ is trivial ($|x|_1 = 1$ for all $x \neq 0$), then by the second step so is $|\cdot|_2$, and they are equal. Otherwise, fix $y \in \mathbb{K}$ with $|y|_1 > 1$; then also $|y|_2 > 1$. Let $x \neq 0$ and write $|x|_1 = |y|_1^{\alpha}$ for some $\alpha \in \mathbb{R}$. For a rational number $\frac{m}{n} > \alpha$ (with $n > 0$) we have $|x|_1 < |y|_1^{m/n}$, so $\left|\frac{x^n}{y^m}\right|_1 < 1$, so $\left|\frac{x^n}{y^m}\right|_2 < 1$, so $|x|_2 < |y|_2^{m/n}$. In the same way, for $\frac{m}{n} < \alpha$ we get $|x|_2 > |y|_2^{m/n}$. Letting $\frac{m}{n} \rightarrow \alpha$ from both sides, $|x|_2 = |y|_2^{\alpha}$. Thus
\[
\frac{\log |x|_1}{\log |y|_1} = \alpha = \frac{\log |x|_2}{\log |y|_2},
\quad\text{so}\quad
\log |x|_1 = s \log |x|_2 \text{ with } s = \frac{\log |y|_1}{\log |y|_2} > 0,
\]
i.e.\ $|x|_1 = |x|_2^s$ for all $x \in \mathbb{K}$.
\end{proof}

\begin{thmbox}{Proposition 2.6}
If $|\cdot|$ is a nontrivial valuation on $\mathbb{Q}$, then the valuation is either equivalent to the absolute value $\|\cdot\|$ or it is equivalent to $|\cdot|_p$ for some prime $p$.
\end{thmbox}

\begin{proof}
Let $|\cdot|$ be an archimedean valuation on $\mathbb{Q}$. I will use $\|\cdot\|$ to represent the absolute value. It is enough to show that $\frac{\log |n|}{\log n}$ is a fixed number for any $n \in \mathbb{N}$, $n \geq 2$. Let $m, n \in \mathbb{N}$ with $m, n \geq 2$. We can write $n$ in base $m$: $n = a_0 + a_1 m + \cdots + a_r m^r$ with $0 \leq a_i < m$, where $r \leq \frac{\log n}{\log m}$. Since $|a_i| = |1 + \cdots + 1| \leq a_i < m$, we get
\[
\begin{aligned}
|n| &\leq \sum_{i=0}^r |a_i| \cdot |m|^i \\
&\leq \sum_{i=0}^r m \cdot \max\{1, |m|\}^r \\
&\leq m\, (r + 1) \max\{1, |m|\}^{\frac{\log n}{\log m}} .
\end{aligned}
\]
If we replace $n$ by $n^k$ (so $r + 1 \leq 1 + k\frac{\log n}{\log m}$), we get
\[
\begin{aligned}
|n|^k &\leq m \left(1 + k\frac{\log n}{\log m}\right) \max\{1, |m|\}^{\frac{k \log n}{\log m}} \\
|n| &\leq m^{\frac1k} \left(1 + k\frac{\log n}{\log m}\right)^{\frac1k} \max\{1, |m|\}^{\frac{\log n}{\log m}} .
\end{aligned}
\]
Letting $k \rightarrow \infty$ we get $|n| \leq \max\{1, |m|\}^{\frac{\log n}{\log m}}$. If $|m| \leq 1$ for some $m \geq 2$, this gives $|n| \leq 1$ for all $n$, so the valuation would be nonarchimedean. Hence $|m| > 1$ for all $m \geq 2$, and $|n| \leq |m|^{\frac{\log n}{\log m}}$. Then if we take the log of both sides we get $\frac{\log |n|}{\log n} \leq \frac{\log |m|}{\log m}$. Since $m, n$ are arbitrary, we can swap them, so $\frac{\log |n|}{\log n}$ is a fixed number $s > 0$ for any $n \in \mathbb{N}$, $n \geq 2$, i.e.\ $|n| = n^s$. For $n \in \mathbb{Z}$ we will have $|n| = |{-1}|\,\|n\|^s = \|n\|^s$, and also $\left|\frac{1}{n}\right| = \frac{1}{|n|} = \left\|\frac1n\right\|^s$. Thus, for $\frac{a}{b} \in \mathbb{Q}$ we get
\[
\left| \frac{a}{b} \right| = |a| \cdot \left| \frac{1}{b} \right| = \|a\|^s \cdot \left\| \frac{1}{b} \right\|^s = \left\| \frac{a}{b} \right\|^s .
\]

Now suppose that the valuation is nonarchimedean, so $|n| \leq 1$ for all $n \in \mathbb{N}$ (Proposition 2.1). Let $p$ be a prime such that $|p| < 1$. If such a prime does not exist, then every prime has $|q| = 1$, and by factorization $|r| = 1$ for all $r \in \mathbb{Q} - \{0\}$: the trivial valuation. So there exists a prime number $p$ with $|p| < 1$. If we look at the set $\mathfrak{a} = \{x \in \mathbb{Z} : |x| < 1\}$, this is an ideal in the ring $\mathbb{Z}$ (by the strong triangle inequality) which contains the ideal $p\mathbb{Z}$; but the latter is a maximal ideal and $1 \notin \mathfrak{a}$, so $\mathfrak{a} = p\mathbb{Z}$. Thus for any prime $q \neq p$ we get $|q| = 1$. Let $|p| = t$ and $r \in \mathbb{Q}$, $r \neq 0$. We write $r = p^m a$ where $p$ divides neither the numerator nor the denominator of $a$, so $|a| = 1$. Then
\[
|r| = |p^m a| = |p|^m = t^m = (p^{-m})^{\frac{-\log t}{\log p}} = |r|_p^s,
\]
where we set $s = \frac{-\log t}{\log p} > 0$.
\end{proof}

\begin{exbox}{Example 2.7}
Let $\mathbb{K}$ be a field. We can define a valuation on the field $\mathbb{K}(x)$ (the field of rational functions) in a similar way to $\mathbb{Q}_p$. Let $p(x)$ be an irreducible monic polynomial in $\mathbb{K}[x]$. To define the exponential valuation on $\mathbb{K}(x)$, let $f(x) \in \mathbb{K}(x)$, $f \neq 0$; then we can write
\[
f(x) = p(x)^m \, \frac{g(x)}{h(x)}, \quad p(x) \nmid g(x)h(x) .
\]
Then $v_p(f) = m$. The multiplicative valuation would be defined as
\[
|f(x)|_p = q^{-v_p(f)} = q^{-m}, \quad\text{where } q \in \mathbb{R} \text{ and } q > 1 .
\]
For $p(x) = x$, the completion of the field $\mathbb{K}(x)$ with respect to the nonarchimedean valuation $|\cdot|_x$ is the field $\mathbb{K}((x))$ of formal Laurent series, which contains the elements in the form
\[
\mathbb{K}((x)) = \left\{ \sum_{i=n}^\infty a_i x^i : n \in \mathbb{Z},\ a_i \in \mathbb{K} \right\}.
\]
The field of formal Laurent series $\mathbb{K}((x))$ is the field of fractions of the ring $\mathbb{K}[[x]]$ of formal power series, which is
\[
\mathbb{K}[[x]] = \left\{ \sum_{i=0}^\infty a_i x^i : a_i \in \mathbb{K} \right\}.
\]
\end{exbox}

Using Hensel's lemma (\textit{we will talk about it later}) we can prove that the algebraic closure of the field of formal Laurent series $\mathbb{K}((x))$, with coefficients from an algebraically closed field $\mathbb{K}$ of characteristic zero, is the field of Puiseux series $\mathbb{K}((x^*))$, which is the union
\[
\bigcup_{n \in \mathbb{N}} \mathbb{K}((x^{\frac1n})) .
\]
Clearly $\mathbb{K}((x^{\frac1n})), \mathbb{K}((x^{\frac1m})) \subset \mathbb{K}((x^{\frac{1}{mn}}))$, so this union is a field. We define the \textit{order} of an element $f(x) \in \mathbb{K}((x^*))$: if $f(x) = \sum_{i=n}^\infty a_i x^{\frac{i}{k}}$ with $a_n \neq 0$, then $\operatorname{ord}(f) = \frac{n}{k}$.

\begin{thmbox}{Proposition 2.8}
If $\mathbb{K}$ is a field with a nonarchimedean valuation $|\cdot|$ and the exponential valuation $v$ (so $|x| = q^{-v(x)}$ for some $q > 1$), then
\begin{itemize}
\item $\mathfrak{o} = \{x \in \mathbb{K} : |x| \leq 1\} = \{x \in \mathbb{K} : v(x) \geq 0\}$ is a ring,
\item $\mathfrak{o}^* = \{x \in \mathbb{K} : |x| = 1\} = \{x \in \mathbb{K} : v(x) = 0\}$ is the group of units of $\mathfrak{o}$,
\item $\mathfrak{p} = \{x \in \mathbb{K} : |x| < 1\} = \{x \in \mathbb{K} : v(x) > 0\}$ is the unique maximal ideal of $\mathfrak{o}$.
\end{itemize}
\end{thmbox}

\begin{proof}
Let $x, y \in \mathfrak{o}$; then $|xy| = |x||y| \leq 1$ and $|x + y| \leq \max\{|x|, |y|\} \leq 1$. Now if $x \in \mathfrak{o}$, then since $|-1| = 1$ we get $|-x| \leq 1$. The other properties follow from the fact that $\mathbb{K}$ is a field, so $\mathfrak{o}$ is a ring. Let $x \in \mathfrak{o}^*$, so $|x| = 1$; then since $1 = |1| = |x\cdot x^{-1}| = |x|\cdot|x^{-1}| = |x^{-1}|$, we get that $x^{-1} \in \mathfrak{o}^*$. Conversely, if $x$ and $x^{-1}$ are both in $\mathfrak{o}$, then $|x| \leq 1$ and $\frac{1}{|x|} \leq 1$, so $|x| = 1$. So $\mathfrak{o}^*$ is exactly the group of units of $\mathfrak{o}$.

Now for the third part: if $x, y \in \mathfrak{p}$ then $|x + y| \leq \max\{|x|, |y|\} < 1$, and if $x \in \mathfrak{o}$ then for any $y \in \mathfrak{p}$ we get that $|xy| = |x||y| < 1$, since $|y| < 1$ and $|x| \leq 1$. So $\mathfrak{p}$ is an ideal. It is the unique maximal ideal: every element of $\mathfrak{o}$ which is not in $\mathfrak{p}$ has $|x| = 1$, so it is a unit; hence a proper ideal contains no element outside $\mathfrak{p}$, i.e.\ every proper ideal is contained in $\mathfrak{p}$.
\end{proof}

We say that a field $\mathbb{K}$ is \textit{complete} with respect to a valuation $|\cdot|$ if every \textit{Cauchy} sequence in $\mathbb{K}$ converges. An integral domain $\mathfrak{o}$ is called a \textit{valuation ring} if for any $x \neq 0$ in $\mathbb{K}$, where $\mathbb{K}$ is its field of fractions, we have either $x \in \mathfrak{o}$ or $x^{-1} \in \mathfrak{o}$. An exponential valuation $v$ is called a \textit{discrete valuation} if there is $s > 0$ such that $v(\mathbb{K}^*) = s\mathbb{Z}$. A discrete valuation is called a \textit{normalized valuation} if $s = 1$. For a normalized valuation, an element $\pi \in \mathfrak{o}$ is called a prime element if $v(\pi) = 1$, and for any $x \in \mathbb{K}^*$ we can write $x = u\pi^m$ where $m \in \mathbb{Z}$ and $u \in \mathfrak{o}^*$, i.e.\ $v(u) = 0$.

\begin{exbox}{Example 2.9}
The completion of the field $\mathbb{Q}$ with respect to the valuation $|\cdot|_p$ is the field of $p$-adic numbers. We can show that in a similar way to showing that the completion of $\mathbb{Q}$ with respect to the absolute value is $\mathbb{R}$.
\end{exbox}

\begin{thmbox}{Proposition 2.10}
If $v$ is a discrete exponential valuation of $\mathbb{K}$, then
\[
\mathfrak{o} = \{ x \in \mathbb{K} : v(x) \geq 0 \}
\]
is a principal ideal domain. Suppose $v$ is normalized. Then the nonzero ideals of $\mathfrak{o}$ are given by $\mathfrak{p}^n = \pi^n \mathfrak{o} = \{ x \in \mathbb{K} : v(x) \geq n \}$, $n \geq 0$, where $\pi$ is a prime element, i.e.\ $v(\pi) = 1$, and one has
\[
\mathfrak{p}^n / \mathfrak{p}^{n+1} \cong \mathfrak{o}/\mathfrak{p} .
\]
\end{thmbox}

\begin{proof}
The valuation is discrete, so there exists $s > 0$ such that $v(\mathbb{K}^*) = s\mathbb{Z}$. Let $I$ be a nonzero ideal in $\mathfrak{o}$. Since the valuation is discrete and $v \geq 0$ on $\mathfrak{o}$, we can define
\[
r = \min_{i \in I,\ i \neq 0} v(i) \in s\mathbb{Z}_{\geq 0},
\]
and let $a \in I$ be an element such that $v(a) = r$. Let $x \in I$; then $v(x) \geq r$. Then the element $xa^{-1} \in \mathfrak{o}$, since $v(xa^{-1}) = v(x) + v(a^{-1}) = v(x) - r \geq 0$. Thus $x = (xa^{-1})\cdot a$, which means $I = (a)$.

Suppose that $v$ is a normalized discrete valuation of $\mathbb{K}$. Let $I$ be a nonzero ideal of $\mathfrak{o}$. From the above it is a principal ideal, so there is $a \in I$ such that $I = (a)$. Now let $\pi \in \mathfrak{o}$ be such that $v(\pi) = 1$. Suppose $v(a) = m$. Then $v(\pi^{-m}a) = 0$, so $u = \pi^{-m}a \in \mathfrak{o}^*$. Thus $a = \pi^m \cdot u$, which means $I = (\pi^m) = \pi^m \mathfrak{o} = \{x : v(x) \geq m\}$.

To show the last result we define the homomorphism (of additive groups) $\phi : \mathfrak{p}^n \rightarrow \mathfrak{o}/\mathfrak{p}$ by $\pi^n x \mapsto x \pmod{\mathfrak{p}}$, $x \in \mathfrak{o}$. It is onto, and an element is in the kernel iff it is of the form $\pi^n \cdot (\pi u) = \pi^{n+1}u$ with $u \in \mathfrak{o}$. Hence the kernel is $\mathfrak{p}^{n+1}$, and that establishes the isomorphism $\mathfrak{p}^n / \mathfrak{p}^{n+1} \cong \mathfrak{o}/\mathfrak{p}$.
\end{proof}

\begin{thmbox}{Hensel's Lemma (second version)}
Let $\mathbb{K}$ be a field which is complete with respect to a nonarchimedean valuation $|\cdot|$, with valuation ring $\mathfrak{o}$ and unique maximal ideal $\mathfrak{p}$. Suppose a primitive polynomial $f(x) \in \mathfrak{o}[x]$ (that is, $\max_i |a_i| = 1$ over its coefficients $a_i$) factors as
\[
f(x) \equiv g_0(x)\, h_0(x) \pmod{\mathfrak{p}},
\]
where $g_0, h_0$ are relatively prime mod $\mathfrak{p}$. Then there exist $g(x)$ and $h(x)$ in $\mathfrak{o}[x]$ such that $f = gh$, $g \equiv g_0$, $h \equiv h_0 \pmod{\mathfrak{p}}$, $\deg(g) = \deg(g_0 \bmod \mathfrak{p})$ and $\deg(h) = \deg(f) - \deg(g)$.
\end{thmbox}

\begin{proof}
We may choose $g_0, h_0$ so that $\deg(g_0) = m$ is the degree of $g_0 \bmod \mathfrak{p}$ (so the leading coefficient of $g_0$ is a unit) and, with $d = \deg(f)$, $\deg(h_0) \leq d - m$. Since $g_0, h_0$ are relatively prime mod $\mathfrak{p}$, there exist $a, b \in \mathfrak{o}[x]$ such that $g_0 a + h_0 b \equiv 1 \pmod{\mathfrak{p}}$. Now we look at the coefficients of the two polynomials
\[
f - g_0 h_0, \qquad g_0 a + h_0 b - 1 .
\]
Then we pick the coefficient with the least valuation $v$ (the largest $|\cdot|$) and call it $\pi$; it is in $\mathfrak{p}$. (If all these coefficients are $0$, then $f = g_0h_0$ and we are done.) We want to find, by induction, $g_1, g_2, \cdots \in \mathfrak{o}[x]$ with $\deg(g_i) < m$, and $h_1, h_2, \cdots \in \mathfrak{o}[x]$ with $\deg(h_i) \leq d - m$, such that
\[
\begin{aligned}
g &= g_0 + g_1 \pi + g_2 \pi^2 + \cdots \\
h &= h_0 + h_1 \pi + h_2 \pi^2 + \cdots
\end{aligned}
\]
and $f \equiv G_n H_n \pmod{\pi^n}$, where
\[
\begin{aligned}
G_n &= g_0 + g_1\pi + \cdots + g_{n-1}\pi^{n-1}, \\
H_n &= h_0 + h_1\pi + \cdots + h_{n-1}\pi^{n-1}.
\end{aligned}
\]
(Here $\equiv \pmod{\pi^n}$ means that the difference has all its coefficients in $\pi^n\mathfrak{o}$.)

For $n = 1$ we have $f \equiv G_1 H_1 = g_0 h_0 \pmod{\pi}$ by the choice of $\pi$. Suppose $f \equiv G_n H_n \pmod{\pi^n}$; we want to find $g_n, h_n \in \mathfrak{o}[x]$ such that
\[
f \equiv (G_n + g_n\pi^n)(H_n + h_n\pi^n) \pmod{\pi^{n+1}} .
\]
Expanding, the terms with $\pi$ exponent more than or equal to $n+1$ will be zero, and $G_n \equiv g_0$, $H_n \equiv h_0 \pmod \pi$:
\[
\begin{aligned}
f &\equiv G_n H_n + g_0 h_n \pi^n + g_n h_0 \pi^n && \pmod{\pi^{n+1}} \\
f - G_n H_n &\equiv \pi^n (g_0 h_n + g_n h_0) && \pmod{\pi^{n+1}} \\
\frac{f - G_n H_n}{\pi^n} &\equiv g_0 h_n + g_n h_0 && \pmod{\pi} .
\end{aligned}
\]
But we also have $g_0 a + h_0 b \equiv 1 \pmod{\pi}$; multiplying by $L(x) = \frac{f - G_n H_n}{\pi^n} \in \mathfrak{o}[x]$ we get
\[
L a g_0 + L b h_0 \equiv L \pmod{\pi} .
\]
Use Euclid division (possible because the leading coefficient of $g_0$ is a unit): $Lb = q(x) g_0(x) + r(x)$ with $\deg(r) < \deg(g_0) = m$. We let $g_n(x) = r(x)$; then we get
\[
L a g_0 + h_0 (q g_0 + g_n) = g_0 (L a + q h_0) + h_0 g_n \equiv L \pmod{\pi} .
\]
Let $h_n$ be $La + qh_0$ with its coefficients in $\pi\mathfrak{o}$ removed. Then $g_0 h_n + h_0 g_n \equiv L \pmod{\pi}$. Moreover $\deg(L) \leq d$ and $\deg(h_0 g_n) < (d - m) + m = d$, so $g_0 h_n \equiv L - h_0 g_n$ has degree at most $d$ mod $\pi$; since the leading coefficient of $g_0$ is a unit, $\deg(h_n) \leq d - m$.

Since $\mathbb{K}$ is complete and $|\pi^n| \rightarrow 0$, the coefficients of $G_n$ and $H_n$ converge, to polynomials $g$ and $h$ with $\deg(g) = m$ and $\deg(h) \leq d - m$, and $f \equiv gh \pmod{\pi^n}$ for every $n$, so $f = gh$. Comparing degrees, $\deg(h) = d - m$.
\end{proof}

\begin{thmbox}{Proposition 2.11}
Let $\mathbb{K}$ be an algebraically closed field with characteristic zero. Then the algebraic closure of the field $\mathbb{K}((x))$ is the Puiseux series field $\mathbb{K}((x^*)) = \bigcup_n \mathbb{K}((x^{\frac1n}))$. The main ideas of the proof are from [3].
\end{thmbox}

\begin{proof}
The field $\mathbb{K}((x))$ is not algebraically closed, because of $F(x, Y) = Y^2 - x$: clearly $x \in \mathbb{K}((x))$ but $\pm x^{\frac12} \notin \mathbb{K}((x))$. Every element of $\mathbb{K}((x^{\frac1k}))$ is algebraic over $\mathbb{K}((x))$, since $\mathbb{K}((x^{\frac1k})) = \mathbb{K}((x))(x^{\frac1k})$ is an extension of degree $k$. Thus if we show that $\mathbb{K}((x^*))$ is algebraically closed, we are done.

It is enough to show that any monic polynomial of degree $n \geq 2$ in $\mathbb{K}((x^*))[Y]$ factors into two polynomials of smaller degree; then, by induction on the degree, every polynomial factors into linear factors. Let
\[
F(x, Y) = Y^n + \left( \sum_{i=r_1}^\infty a_{1i} x^{\frac{i}{k_1}} \right) Y^{n-1} + \cdots + \left( \sum_{i=r_n}^\infty a_{ni} x^{\frac{i}{k_n}} \right), \quad a_{ji} \in \mathbb{K} .
\]
Using the Tschirnhausen transformation $Y = Y' - \frac{\text{coefficient of } Y^{n-1}}{n}$ (possible since the characteristic is zero) we can assume that the coefficient of $Y^{n-1}$ is zero. We can also assume that the constant term is not zero (otherwise $Y$ is a factor). Thus we have
\[
F(x, Y) = Y^n + \left( \sum_{i=r_2}^\infty a_{2i} x^{\frac{i}{k_2}} \right) Y^{n-2} + \cdots + \left( \sum_{i=r_n}^\infty a_{ni} x^{\frac{i}{k_n}} \right),
\]
where, for each coefficient $c_j$ which is not zero, $a_{jr_j} \neq 0$, so $\operatorname{ord}(c_j) = \frac{r_j}{k_j}$. Let $r = \min_j \frac{r_j}{j k_j}$ over the nonzero coefficients $c_j$. Using a suitable substitution $z^m = x$, for example $m = \operatorname{LCM}_{j}(j k_j)$, let $m_j = \frac{m}{k_j}$; then $mr$ is an integer and
\[
F(z, Y) = Y^n + \left( \sum_{i=r_2}^\infty a_{2i} z^{i m_2} \right) Y^{n-2} + \cdots + \left( \sum_{i=r_n}^\infty a_{ni} z^{i m_n} \right) .
\]
Set $Y = z^{mr} U$. Then we get
\[
\begin{aligned}
F(z, z^{mr}U) &= U^n z^{nmr} + \left( \sum_{i=r_2}^\infty a_{2i} z^{i m_2} \right) z^{(n-2)mr} U^{n-2} + \cdots + \left( \sum_{i=r_n}^\infty a_{ni} z^{i m_n} \right) \\
&= z^{nmr} \left[ U^n + \left( \sum_{i=r_2}^\infty a_{2i} z^{i m_2 - 2mr} \right) U^{n-2} + \cdots + \left( \sum_{i=r_n}^\infty a_{ni} z^{i m_n - nmr} \right) \right] .
\end{aligned}
\]
The exponent of $z$ in each coefficient is non-negative: the smallest exponent is $r_j m_j - jmr \geq r_j \frac{m}{k_j} - jm \frac{r_j}{j k_j} = 0$, with equality for the $j$ where the minimum $r$ is attained. Now consider the polynomial
\[
G(z, U) = U^n + \left( \sum_{i=r_2}^\infty a_{2i} z^{i m_2 - 2mr} \right) U^{n-2} + \cdots + \left( \sum_{i=r_n}^\infty a_{ni} z^{i m_n - nmr} \right) \in \mathbb{K}[[z]][U] .
\]
If we look at $G(0, U)$, this polynomial factors into two relatively prime polynomials, because $\mathbb{K}$ is algebraically closed and the coefficient of $U^{n-1}$ is zero. Indeed, $G(0, U)$ is not $U^n$, since the coefficient where the minimum $r$ is attained has a nonzero constant term. If $G(0, U)$ had only one root $\alpha$, then $G(0, U) = (U - \alpha)^n$, whose coefficient of $U^{n-1}$ is $-n\alpha = 0$, so $\alpha = 0$ (characteristic zero) and $G(0, U) = U^n$, a contradiction. So $G(0, U)$ has two different roots $\alpha \neq \beta$, and $G(0, U) = (U - \alpha)^e\, h_0(U)$ with $h_0(\alpha) \neq 0$ and $1 \leq e < n$: two relatively prime factors. Hence, using Hensel's lemma (second version) for the complete field $\mathbb{K}((z))$, with valuation ring $\mathbb{K}[[z]]$ and $\mathbb{K}[[z]]/(z) = \mathbb{K}$, we get $G = gh$ with $\deg(g) = e$ and $\deg(h) = n - e$. Then $F(z, Y) = z^{nmr}\, G(z, z^{-mr}Y)$ factors into two polynomials of degrees $e$ and $n - e$ over $\mathbb{K}((z)) = \mathbb{K}((x^{\frac1m})) \subset \mathbb{K}((x^*))$, which is the desired.
\end{proof}

\begin{thmbox}{Ostrowski's Theorem}
\index{Ostrowski Theorem}Let $\mathbb{K}$ be a field of characteristic $0$. Suppose that it is complete with respect to an archimedean valuation $|\cdot|$. Then $\mathbb{K}$ is isomorphic to either $\mathbb{R}$ or $\mathbb{C}$.
\end{thmbox}

\begin{proof}
First, I will show that $\mathbb{K}$ contains a subfield isomorphic to $\mathbb{R}$. Let $1_{\mathbb{K}}$ be the multiplicative identity of $\mathbb{K}$; then we consider the ring homomorphism
\[
\phi : \mathbb{Z} \rightarrow \mathbb{K}, \quad \phi(1) = 1_{\mathbb{K}},
\]
and $\phi(n) = n\cdot 1_{\mathbb{K}}$. It is one to one because the characteristic is $0$. Thus $\mathbb{K}$ contains a copy of $\mathbb{Z}$. Now since $\underbrace{1_{\mathbb{K}} + 1_{\mathbb{K}} + \cdots + 1_{\mathbb{K}}}_{n \text{ times}} = n 1_{\mathbb{K}} \in \mathbb{K}$ is not zero, its multiplicative inverse is contained in $\mathbb{K}$, and it can be represented as $\frac1n 1_{\mathbb{K}}$. The homomorphism $\phi$ can be extended to $\mathbb{Q}$: we map $\frac1n$ to $\frac1n\cdot 1_{\mathbb{K}}$, so by that $\mathbb{K}$ contains a copy of $\mathbb{Q}$, and
\[
\phi : \mathbb{Q} \rightarrow \mathbb{K}
\]
will be a homomorphism between fields. Now to show that $\mathbb{K}$ contains a copy of $\mathbb{R}$, we define a valuation on $\mathbb{Q}$, call it $|\cdot|_{\mathbb{Q}}$, as follows:
\[
|\cdot|_{\mathbb{Q}} : \mathbb{Q} \rightarrow \mathbb{R}, \quad |r|_{\mathbb{Q}} = |\phi(r)|_{\mathbb{K}} .
\]
This defines an archimedean valuation on $\mathbb{Q}$ (the values $|n\cdot 1_{\mathbb{K}}|$ are not bounded). But we proved earlier (Proposition 2.6) that any archimedean valuation on $\mathbb{Q}$ is equivalent to $\|\cdot\|$, the absolute value; thus $|\cdot|_{\mathbb{Q}} = \|\cdot\|^s$ for some $s \in \mathbb{R}^+$. Now let $v \in \mathbb{R}$; then there exists a sequence of rational numbers $r_n \rightarrow v$ with respect to the absolute value. Then $\phi(r_n)$ is a Cauchy sequence in $\mathbb{K}$, since $|\phi(r_n) - \phi(r_m)|_{\mathbb{K}} = |r_n - r_m|_{\mathbb{Q}} = \|r_n - r_m\|^s$, so it converges, since $\mathbb{K}$ is complete, and we set $\phi(v) = \lim \phi(r_n)$ (the limit does not depend on the sequence). Hence $\mathbb{K}$ contains a copy of $\mathbb{R}$, on which $|\cdot|_{\mathbb{K}} = \|\cdot\|^s$. Without loss of generality we can assume $\mathbb{R} \subset \mathbb{K}$, with $|x| = \|x\|^s$ for $x \in \mathbb{R}$.

Now we will show that any $\xi \in \mathbb{K}$ satisfies a quadratic polynomial in $\mathbb{R}[x]$. Consider the following function $f : \mathbb{C} \rightarrow \mathbb{R}$:
\[
f(z) = |\xi^2 - (z + \overline{z})\xi + z\overline{z}|_{\mathbb{K}} .
\]
Note that $z + \overline{z}$ and $z\overline{z}$ are real numbers, so the expression is in $\mathbb{K}$. The function $f(z)$ is continuous: the maps $z \mapsto z + \overline{z}$ and $z \mapsto z\overline{z}$ are continuous from $\mathbb{C}$ to $\mathbb{R}$, and $f$ is obtained from them by the operations of $\mathbb{K}$ and the valuation, which are continuous (for example $\big|\,|a|_{\mathbb{K}} - |b|_{\mathbb{K}}\,\big| \leq |a - b|_{\mathbb{K}}$). Also $f(z) \geq \|z\|^{2s} - \|z + \overline{z}\|^s\,|\xi| - |\xi|^2 \rightarrow \infty$ as $\|z\| \rightarrow \infty$. Hence $f(z)$ has a minimum $m$; define
\[
S = \{ z \in \mathbb{C} : f(z) = m \}.
\]
If we prove that $m = 0$, then we will have $z \in \mathbb{C}$ such that $f(z) = 0$, i.e.\ $\xi^2 - (z + \overline{z})\xi + z\overline{z} = 0$. Suppose, on the contrary, that $m > 0$. The set $S$ is nonempty, bounded and closed. It is bounded because $\lim_{\|z\| \rightarrow \infty} f(z) = \infty$, and it is closed because it is the inverse image of a closed set. Now if we look at $\{ \|z\| : z \in S \}$, this is a closed bounded subset of $\mathbb{R}$, so it has a maximum, attained at some $z_0 \in S$. Let $\epsilon \in \mathbb{R}$ be such that $\epsilon > 0$ and $|\epsilon| = \epsilon^s < m$, and define the function
\[
h(x) = x^2 - (z_0 + \overline{z}_0)x + z_0\overline{z}_0 + \epsilon .
\]
The function $h(x)$ has no real roots, because it has a negative discriminant: $(z_0 + \overline{z}_0)^2 - 4(z_0\overline{z}_0 + \epsilon) = (z_0 - \overline{z}_0)^2 - 4\epsilon = -4(\operatorname{Im} z_0)^2 - 4\epsilon < 0$. Let $z_1, \overline{z}_1 \in \mathbb{C}$ be the two complex roots of $h(x)$. Since the constant term of a quadratic polynomial is the product of its roots, we get that $z_1\overline{z}_1 = z_0\overline{z}_0 + \epsilon$. Thus $\|z_1\| > \|z_0\|$, so $z_1 \notin S$, i.e.\ $f(z_1) > m$.

Define the function
\[
G(x) = (h(x) - \epsilon)^n - (-\epsilon)^n .
\]
Then $G(z_1) = (h(z_1) - \epsilon)^n - (-\epsilon)^n = (-\epsilon)^n - (-\epsilon)^n = 0$. $G(x)$ is a real polynomial of degree $2n$; suppose $G(x)$ has the roots $\alpha_1, \cdots, \alpha_{2n}$ and assume $z_1 = \alpha_1$. Then $G(x)$ will factor completely as
\[
G(x) = \prod_{i=1}^{2n} (x - \alpha_i) .
\]
But since $G$ has real coefficients, we also have $G(x) = \prod (x - \overline{\alpha}_i)$. Thus
\[
(G(x))^2 = \prod (x - \alpha_i) \prod (x - \overline{\alpha}_i) = \prod \big( x^2 - (\alpha_i + \overline{\alpha}_i)x + \alpha_i\overline{\alpha}_i \big) .
\]
Since $f(z) = |\xi^2 - (z + \overline{z})\xi + z\overline{z}| \geq m$ for any $z \in \mathbb{C}$, we get that
\[
|G(\xi)|^2 = \prod_{i=1}^{2n} f(\alpha_i) \geq f(z_1)\, m^{2n-1} .
\]
Also we have, since $h(\xi) - \epsilon = \xi^2 - (z_0 + \overline{z}_0)\xi + z_0\overline{z}_0$,
\[
|G(\xi)| = |(h(\xi) - \epsilon)^n - (-\epsilon)^n| \leq |h(\xi) - \epsilon|^n + |\epsilon|^n = (f(z_0))^n + |\epsilon|^n = m^n + |\epsilon|^n .
\]
Combining the two inequalities we get
\[
\begin{aligned}
f(z_1)\, m^{2n-1} &\leq \left( m^n + |\epsilon|^n \right)^2 \\
\frac{f(z_1)}{m} &\leq \left( 1 + \left( \frac{|\epsilon|}{m} \right)^n \right)^2 .
\end{aligned}
\]
Note that $0 < |\epsilon| < m$, so $\frac{|\epsilon|}{m} < 1$. Thus if we take the limit as $n \rightarrow \infty$ we get
\[
f(z_1) \leq m .
\]
But then $z_1 \in S$, which contradicts the fact that $\|z_1\| > \|z_0\|$. So $m = 0$.

Finally, every $\xi \in \mathbb{K}$ is a root of a real polynomial of degree $2$. If $\mathbb{K} \neq \mathbb{R}$, take $\xi \in \mathbb{K} - \mathbb{R}$; then $\mathbb{R}(\xi)$ is an extension of degree $2$ of $\mathbb{R}$, so $\mathbb{R}(\xi) \cong \mathbb{C}$, and $\mathbb{K}$ contains a copy of $\mathbb{C}$. Every element of $\mathbb{K}$ is algebraic over this copy of $\mathbb{C}$, which is algebraically closed, so $\mathbb{K} \cong \mathbb{C}$. Otherwise $\mathbb{K} = \mathbb{R}$.
\end{proof}

\section{Extension of valuations}

If we take a field $\mathbb{K}$ which is complete with respect to a valuation $|\cdot|$, archimedean or nonarchimedean, and consider an algebraic extension $\mathbb{L}|\mathbb{K}$, we can extend the valuation $|\cdot|$ to $\mathbb{L}$ uniquely. In the following example we will see the case of a quadratic extension.

\begin{exbox}{Example 3.1}
Let $\mathbb{K}$ be a complete field with valuation $|\cdot|_{\mathbb{K}}$, of characteristic not $2$. Suppose $p(x) = x^2 + bx + c$ is an irreducible polynomial in $\mathbb{K}[x]$. Let $\theta = \sqrt{b^2 - 4c}$; then the polynomial factors completely in $\mathbb{K}(\theta)[x]$ as
\[
x^2 + bx + c = \left( x - \frac{-b - \theta}{2} \right)\left( x - \frac{-b + \theta}{2} \right) .
\]
Let $\mathbb{L} = \mathbb{K}(\theta) = \{ \alpha + \theta\beta : \alpha, \beta \in \mathbb{K} \}$. We want to extend the valuation $|\cdot|_{\mathbb{K}}$ to the extension $\mathbb{L}$ as $|\cdot|_{\mathbb{L}}$, and we want the extension of the valuation to restrict to the field $\mathbb{K}$, i.e.\ for any $\gamma \in \mathbb{K}$, $|\gamma|_{\mathbb{K}} = |\gamma|_{\mathbb{L}}$. If $\alpha + \theta\beta \in \mathbb{L}$, then $(\alpha + \theta\beta)(\alpha - \theta\beta) = \alpha^2 - (b^2 - 4c)\beta^2 \in \mathbb{K}$, so we have
\begin{equation}\label{eq:val-ext}
|(\alpha + \theta\beta)(\alpha - \theta\beta)|_{\mathbb{K}} = |\alpha + \theta\beta|_{\mathbb{L}}\, |\alpha - \theta\beta|_{\mathbb{L}} .
\end{equation}
If $\alpha = 0$, $\beta = 1$, we get $|\theta|_{\mathbb{L}}\,|-\theta|_{\mathbb{L}} = |\theta|_{\mathbb{L}}^2 = |-(b^2 - 4c)|_{\mathbb{K}}$, so $|\theta|_{\mathbb{L}} = \sqrt{|b^2 - 4c|_{\mathbb{K}}}$. Note that $|\theta|_{\mathbb{L}}^2 = |b^2 - 4c|_{\mathbb{K}}$ is the valuation of the constant term in the minimal polynomial of $\theta$, which is $x^2 - (b^2 - 4c)$.
\end{exbox}

\begin{defnbox}{Definition}
Let $\mathbb{L}|\mathbb{K}$ be an algebraic extension. If $\alpha \in \mathbb{L}$, then it is algebraic over $\mathbb{K}$. Suppose $[\mathbb{K}(\alpha) : \mathbb{K}] = n$, with minimal polynomial
\[
p(x) = x^n + c_1 x^{n-1} + \cdots + c_n, \quad\text{and } p(\alpha) = 0 .
\]
We define the \textit{norm} of $\alpha$ as $N_{\mathbb{K}(\alpha)|\mathbb{K}}(\alpha) = \det(A_\alpha)$, where $A_\alpha$ is the linear transformation $A_\alpha : \mathbb{K}(\alpha) \rightarrow \mathbb{K}(\alpha)$, $A_\alpha(x) = \alpha x$, considering $\mathbb{K}(\alpha)$ as a vector space over $\mathbb{K}$ with basis $\{1, \alpha, \alpha^2, \cdots, \alpha^{n-1}\}$. Let $\alpha^i = (0, \cdots, 1, 0, \cdots, 0)^T$ be the vector with $1$ in the $(i+1)$-th coordinate and the remaining zeroes. Then $A_\alpha$ will send $\alpha^i \mapsto \alpha^{i+1}$, except $\alpha^{n-1} \mapsto \alpha^n = -c_n - \cdots - c_1\alpha^{n-1}$. Thus
\[
A_\alpha = \begin{pmatrix}
0 & 0 & \cdots & 0 & -c_n \\
1 & 0 & \cdots & 0 & -c_{n-1} \\
0 & 1 & \cdots & 0 & -c_{n-2} \\
\vdots & \vdots & \ddots & \vdots & \vdots \\
0 & 0 & \cdots & 1 & -c_1
\end{pmatrix},
\qquad N_{\mathbb{K}(\alpha)|\mathbb{K}}(\alpha) = \det(A_\alpha) = (-1)^n c_n .
\]
\end{defnbox}

\begin{thmbox}{Theorem 3.2}
Let $\mathbb{K}$ be a complete field with respect to the valuation $|\cdot|$. Suppose $\mathbb{L}|\mathbb{K}$ is an algebraic extension of $\mathbb{K}$. Then the valuation $|\cdot|$ extends uniquely to $\mathbb{L}$, by the formula
\[
|\alpha|_{\mathbb{L}} = \sqrt[n]{|N_{\mathbb{K}(\alpha)|\mathbb{K}}(\alpha)|_{\mathbb{K}}},
\]
where $n = [\mathbb{K}(\alpha) : \mathbb{K}]$.
\end{thmbox}

\begin{proof}
Suppose the valuation $|\cdot|_{\mathbb{K}}$ is nonarchimedean, and suppose $\mathbb{L}|\mathbb{K}$ is an algebraic extension. Let $\mathfrak{o}$ be the valuation ring of $\mathbb{K}$. First, we need two facts.

(a) \emph{If $g(x) = a_d x^d + \cdots + a_0 \in \mathbb{K}[x]$ is irreducible with $a_0 a_d \neq 0$, then $\max_i |a_i| = \max\{|a_0|, |a_d|\}$.} Multiplying by a constant, we can assume $\max_i |a_i| = 1$, i.e.\ $g$ is primitive. If $\max\{|a_0|, |a_d|\} < 1$, let $r$ be the smallest index with $|a_r| = 1$; then $0 < r < d$ and
\[
g(x) \equiv x^r (a_r + a_{r+1}x + \cdots + a_d x^{d-r}) \pmod{\mathfrak{p}},
\]
two relatively prime factors (the second has nonzero constant term $a_r$ mod $\mathfrak{p}$). By Hensel's lemma (second version) $g$ is a product of polynomials of degrees $r$ and $d - r$, which contradicts that $g$ is irreducible.

(b) \emph{For $\alpha \in \mathbb{L}$: $|N_{\mathbb{K}(\alpha)|\mathbb{K}}(\alpha)| \leq 1$ if and only if the minimal polynomial of $\alpha$ has all its coefficients in $\mathfrak{o}$.} Indeed, by (a) applied to the minimal polynomial $p(x)$ (with $a_d = 1$), its coefficients are all in $\mathfrak{o}$ iff $|c_n| \leq 1$, and $|c_n| = |N_{\mathbb{K}(\alpha)|\mathbb{K}}(\alpha)|$.

Now we show that $|\alpha|_{\mathbb{L}} = |N_{\mathbb{K}(\alpha)|\mathbb{K}}(\alpha)|^{1/n}$ defines a nonarchimedean valuation on $\mathbb{L}$ which extends $|\cdot|$. Suppose $\alpha, \beta \in \mathbb{L}$.
\begin{itemize}
\item If $\alpha \in \mathbb{K}$ then $n = 1$ and $N(\alpha) = \alpha$, so $|\alpha|_{\mathbb{L}} = |\alpha|$: it extends $|\cdot|$.
\item $|\alpha|_{\mathbb{L}} = 0$ iff $N(\alpha) = 0$ iff $c_n = 0$ iff $\alpha = 0$ (the minimal polynomial is irreducible).
\item $|\alpha\beta|_{\mathbb{L}} = |\alpha|_{\mathbb{L}} |\beta|_{\mathbb{L}}$: let $\mathbb{M} = \mathbb{K}(\alpha, \beta)$, a finite extension. For $\gamma \in \mathbb{M}$, the matrix of multiplication by $\gamma$ on $\mathbb{M}$ is made of $[\mathbb{M} : \mathbb{K}(\gamma)]$ copies of $A_\gamma$ (in a basis $\{\mu_j \gamma^i\}$ with $\{\mu_j\}$ a basis of $\mathbb{M}$ over $\mathbb{K}(\gamma)$), so $N_{\mathbb{M}|\mathbb{K}}(\gamma) = N_{\mathbb{K}(\gamma)|\mathbb{K}}(\gamma)^{[\mathbb{M} : \mathbb{K}(\gamma)]}$ and
\[
|\gamma|_{\mathbb{L}} = |N_{\mathbb{M}|\mathbb{K}}(\gamma)|^{1/[\mathbb{M} : \mathbb{K}]} .
\]
Since $N_{\mathbb{M}|\mathbb{K}}(\alpha\beta) = \det(A_{\alpha\beta}) = \det(A_\alpha A_\beta) = N_{\mathbb{M}|\mathbb{K}}(\alpha)\,N_{\mathbb{M}|\mathbb{K}}(\beta)$ (multiplication maps on $\mathbb{M}$), the valuation is multiplicative.
\item $|\alpha + \beta|_{\mathbb{L}} \leq \max\{|\alpha|_{\mathbb{L}}, |\beta|_{\mathbb{L}}\}$: we can suppose $|\alpha|_{\mathbb{L}} \leq |\beta|_{\mathbb{L}}$, $\beta \neq 0$, and divide by $\beta$; so we have to show that $|\gamma|_{\mathbb{L}} \leq 1 \Rightarrow |\gamma + 1|_{\mathbb{L}} \leq 1$. If $|\gamma|_{\mathbb{L}} \leq 1$, then by (b) the minimal polynomial $p(x)$ of $\gamma$ has its coefficients in $\mathfrak{o}$. The minimal polynomial of $\gamma + 1$ is $p(x - 1)$, which also has its coefficients in $\mathfrak{o}$, so by (b) again $|\gamma + 1|_{\mathbb{L}} \leq 1$.
\end{itemize}

Uniqueness: suppose $|\cdot|'$ is another valuation on $\mathbb{L}$ extending $|\cdot|$. First, if $|\alpha|_{\mathbb{L}} \leq 1$ then $|\alpha|' \leq 1$. Indeed, $\alpha$ satisfies $\alpha^n + c_1\alpha^{n-1} + \cdots + c_n = 0$ with $c_i \in \mathfrak{o}$ by (b). If $|\alpha|' > 1$, then dividing by $\alpha^n$,
\[
1 = -c_1\alpha^{-1} - \cdots - c_n\alpha^{-n},
\]
and the right side has $|\cdot|' < 1$, since $|c_i|' = |c_i| \leq 1$ and $|\alpha^{-1}|' < 1$ (we use that $|\cdot|'$ is nonarchimedean, which holds because it is bounded on $\mathbb{N} \subset \mathbb{K}$ by Proposition 2.1). This is a contradiction. Now let $\alpha \neq 0$, $n = [\mathbb{K}(\alpha) : \mathbb{K}]$ and $a = N_{\mathbb{K}(\alpha)|\mathbb{K}}(\alpha) \in \mathbb{K}$. Then $|\alpha^n / a|_{\mathbb{L}} = |\alpha|_{\mathbb{L}}^n / |a| = 1$, so $|\alpha^n / a|' \leq 1$, i.e.\ $(|\alpha|')^n \leq |a|' = |a| = |\alpha|_{\mathbb{L}}^n$, so $|\alpha|' \leq |\alpha|_{\mathbb{L}}$. Doing the same for $\alpha^{-1}$ gives $|\alpha|' \geq |\alpha|_{\mathbb{L}}$. So $|\cdot|' = |\cdot|_{\mathbb{L}}$.

Finally, if $|\cdot|$ is archimedean, then by Ostrowski's theorem $\mathbb{K}$ is $\mathbb{R}$ or $\mathbb{C}$, with $|x| = \|x\|^s$ on $\mathbb{R}$ (as in its proof). Here $s \leq 1$, since $2^s = |1 + 1| \leq |1| + |1| = 2$. If $\mathbb{K} = \mathbb{C}$ there is nothing to prove, because $\mathbb{C}$ is algebraically closed, so $\mathbb{L} = \mathbb{C}$. If $\mathbb{K} = \mathbb{R}$, the only algebraic extension which is not trivial is $\mathbb{C}|\mathbb{R}$. There, for $z \notin \mathbb{R}$, the minimal polynomial of $z$ is $x^2 - (z + \overline{z})x + z\overline{z}$, so $N_{\mathbb{C}|\mathbb{R}}(z) = z\overline{z}$, and the formula gives $|z|_{\mathbb{C}} = \sqrt{\|z\overline{z}\|^s} = \|z\|^s$ (also for $z \in \mathbb{R}$). This is a valuation: the only property which is not clear is the triangle inequality, and it holds because $\|z + w\|^s \leq (\|z\| + \|w\|)^s \leq \|z\|^s + \|w\|^s$ for $0 < s \leq 1$.

For uniqueness, let $|\cdot|'$ be any valuation on $\mathbb{C}$ extending $|\cdot|$. Since $(|i|')^2 = |{-1}| = 1$, we have $|i|' = 1$. For $z = a + bi$ with $a, b \in \mathbb{R}$,
\[
|z|' \leq |a| + |b|\,|i|' = \|a\|^s + \|b\|^s \leq 2\|z\|^s = 2|z|_{\mathbb{C}} .
\]
Applying this to $z^k$ gives $(|z|')^k \leq 2|z|_{\mathbb{C}}^k$, i.e.\ $|z|' \leq 2^{\frac1k}|z|_{\mathbb{C}}$, and letting $k \rightarrow \infty$, $|z|' \leq |z|_{\mathbb{C}}$. Applying the same to $z^{-1}$ (for $z \neq 0$) gives $|z|' \geq |z|_{\mathbb{C}}$. So $|\cdot|' = |\cdot|_{\mathbb{C}}$.
\end{proof}

\section*{References}

\begin{enumerate}[label={[\arabic*]}, leftmargin=*]
\item J.~Neukirch, \textit{Algebraic Number Theory}, Grundlehren der mathematischen Wissenschaften 322, Springer, 1999.
\item W.~Stein, \textit{Algebraic Number Theory, a Computational Approach}, 2012.
\item K.~J.~Nowak, Some elementary proofs of Puiseux's theorems, \textit{Universitatis Iagellonicae Acta Mathematica} 38 (2000), 279--282.
\end{enumerate}

\section*{Appendix: the $p$-adics Python codes}

\begin{rawhtml}
<span id="programs"></span>
\end{rawhtml}

The first code converts rational numbers into $p$-adics. The second code checks if a number has a square root in the $p$-adics.

\begin{webonly}
Both programs run below, in your browser: change the numbers and press Run. They follow the same steps as the Python code, with exact integers, so long expansions are correct. Under each one you can read the Python code and download it.
\end{webonly}

\begin{rawhtml}
<div class="padic-tool" data-tool="rational" data-src="../../latex/projects/valuations-code/padic_tool.py"></div>
<div class="padic-tool" data-tool="sqrt" data-src="../../latex/projects/valuations-code/quadratic.py"></div>
<script src="padic-tools.js"></script>
\end{rawhtml}

\begin{pdfonly}
\paragraph{The first code} (the file \texttt{padic\_tool.py}):
\lstinputlisting[language=Python]{valuations-code/padic_tool.py}

\paragraph{The second code} (the file \texttt{quadratic.py}):
\lstinputlisting[language=Python]{valuations-code/quadratic.py}
\end{pdfonly}

\end{document}
